\documentclass{jsarticle}
\usepackage{amsmath}
\usepackage{amssymb}
\begin{document}
\title{}
\author{}
\date{}
\maketitle

   習い事は、成長したときに、自分の取り柄になるのを、
   2015，2016年の剣道と居合道で、親に感謝することの大切さを
   最近、ふとおもうことです。小学校時代は、全然習うのが嬉しくなかったが、
   何年経っても、技法は廃れていないのが嬉しかった。
   彩さんも、子どもに習い事をSRS速読法をさせているのは、
   自分も便乗したい気持ちで有り、いいとおもいます。
   速読法は、忍耐力は疑問ですが、協調性はあるとおもい、
   子どもがまだ幼いのに、家の手伝いをするすごい親思いの良いことおもいます。
   彩さんが怒ったら怖いのか、「いつも優しくしてくれて、ありがとうございます。」
   を言っている、美穂さんは、野口さんと同じなのかなと、彩さん少し心配です。
   かかあ天下はだめですよ。亭主関白もだめですよ。
   北川悠仁さんに、あそこを、彩子さんと美穂さんが生まれるときに、見せている
   という、北川悠仁さんかわいそうです。人の言えた義理でもないのは、
   承知の上ですが、大和撫子やってくうれていたら、今の論文を書いている手柄を
   彩さんに軍配を上げたかもしれなかったかも。家のネコは、食べ物が良かったの
   ですよと、ボディーランゲージで話しています。いつか、どこかで、
   逢いましょう。

\end{document}
